\documentclass{article}

\usepackage{fancyvrb}
\usepackage{xcolor}
\usepackage{listings}
\usepackage{lipsum}
\usepackage{minted}

\definecolor{pythonback}{rgb}{0.95,0.95,0.92}
\definecolor{pythoncomment}{rgb}{0,0.6,0}
\definecolor{numbercolor}{rgb}{0.5,0.5,0.5}
\definecolor{stringcolor}{rgb}{0.58,0,0.82}

\lstdefinestyle{PythonCode}{
	backgroundcolor = \color{pythonback},
	commentstyle = \color{pythoncomment},
	numberstyle = \color{numbercolor},
	stringstyle = \color{stringcolor},
	keywordstyle = \color{blue},
	basicstyle = \ttfamily\small,
	breaklines = true,
	captionpos = b,
	numbers = left,
	numbersep = 20pt,
	frame = single,
	framerule = 1pt,
	tabsize = 2,
	showstringspaces=false,
	escapechar = \%,
	framesep = \baselineskip
}

\usemintedstyle[python]{igor}

% \surroundwithmdframed{minted}
% \tcbuselibrary{minted}, then use tcblisting env.

\setminted[python]{
   linenos,
   mathescape,
   baselinestretch=1.5,
   bgcolor=gray!10,
   breaklines,
   escapeinside=||,
   frame=lines,
   framerule=1.5pt,
   framesep=10pt,
   highlightlines={7,8},
   highlightcolor=cyan!10,
   tabsize=4,
   texcl=true
}

\begin{document}
	{\noindent \Huge \textbf{Code Listing}}
	
	\vspace{2cm}
	
	\section{The Basic Verbatim Environment}
	\begin{verbatim}
		for i in range(len(X)):
   			# iterate through columns
   			for j in range(len(X[0])):
       			result[j][i] = X[i][j]
	\end{verbatim}
	
	This is an inline example of \verb|\verbatim| in \LaTeX.
	
	\begin{Verbatim}[frame=single,framerule=1mm,
					 framesep=5mm,
					 fillcolor=\color{gray!50}]
	An enhanced verbatim text using fancyvrb package.
	\end{Verbatim}
	
	\vspace{1cm}
	
	\section{The listings Package}
	\subsection{List Inline}
	This is an inline Example of listings 						package \fbox{\lstinline&for i in 						range(len(X)):&} -- and that was the code.
	
	\subsection{Normal Blocks of Code: lstlisting}
	\begin{lstlisting}
		def my_function():
  			print("Hello from a function")
	\end{lstlisting}
	
	\subsection{Code From Source File}
	\lstinputlisting{../Listed-Code-Source-Files/source_file.java}
	
	\subsection{Package Control Options}
	\begin{lstlisting}[style=PythonCode,
					   name=PyBlock,
					   language=Python,
					   frame = ltr,
					   caption=Geometric Series Sum.]
	# Question No. 1
	def sum_geometric_series(a,r,n):
    	term = 0
    	# Start Numbering at One for Right Exponent
    	for i in range(1, n+1): 
        	term += a * pow(r, i) # Accumulating
    	series_sum = a + term
    	return series_sum
	
	# Example
	print("The Result is:",sum_geometric_series(2,3,4))
	%\phantom{-----------}%
	\end{lstlisting}
	
	\bigskip
	\lipsum[1]
	\bigskip
	
	% Emphsizing a function name
	\lstset{emph = approximate_e, 									emphstyle = \color{brown}}
	\begin{lstlisting}[style=PythonCode,
					   name=PyBlock,
					   language=Python,
					   frame = lbr,
					 caption=Approximate Euler Number.]
	%\phantom{-----------}%
	# Question No. 2
	def approximate_e(n):
   		e = pow((1+1/2), n )
    	return e
	
	# Example
	print(approximate_e(2))
	\end{lstlisting}
	
	\pagebreak
	
	\subsection{Listing Specific Range of Lines}			\lstset{showlines=true}
	\lstinputlisting[language=Python,
					 style=PythonCode,
					 linerange={13-18}]
					 {../Listed-Code-Source-Files/input_range-example.py}
	%%%%% The sixth empty line: 
		% By default the package ignores any blank lines in the code but the option "showlines" when set to true allows for printing.
					 
	
	\section{The minted Package}
	\begin{minted}{python}
	# \underline{Question} No. 8: $F = [F_x, F_y]$
	import numpy as np
	from math import *
	from numpy import |\colorbox{green}{array}|
	
	def force_vec(F1, F2, theta1, theta2):
	    Rx = F1 * sin(radians(theta1)) +  F2 * cos(radians(theta2))
	    Ry = F1 * cos(radians(theta1)) +  F2 * sin(radians(theta2))
	    R = [Rx, Ry]
	    return R
	
	# Example
	F1 = 250 
	F2 = 450
	theta1 = 30
	theta2 = 35
	print("Result:", force_vec(F1, F2, theta1, theta2))
	\end{minted}
	
	% \inputminted[firstline=, lastline=]{lang}{PATH}
	
	\bigskip
	%% However: \mint{lang}|CODE| "OneLine"
	\mint{python}|import library|
	
	%% For Inline: \mintinline{lang}{CODE}
	This is an Inline Code \fbox{\mintinline{python}{Print("Hello World !!")}} Offered by minted Package.
	
\end{document}